\scanpage{313}

\mspage{52}
% archive sheet number: 312

Se donner $f\in\operatorname{Fl}(\underline{\mathrm{Ty}})$ revient à se donner une suite ordonnée de générateurs $f_{\varepsilon\,\unsure{;|,}\ell}$, donc une suite
\[
\begin{aligned}
&\deleted{\unsure{f_{(\varepsilon_1,\ell_1)}|f_{(\varepsilon_i,\ell_i)}}}\\
&f_{\varepsilon_1,\ell_1},\quad f_{\varepsilon_2,\ell_2}\quad\cdots\quad f_{\varepsilon_n,\ell_n}\\
&\big((\varepsilon_i,\ell_i)\big)_{1\leq i\leq n}\in[\{\pm1\}\times\mathbb N]^n.
\end{aligned}
\]

On a alors
\[
\begin{aligned}
x&=\Delta_{\varepsilon_1}^{\deleted{\unsure{\ell_1}}}\Delta_{\varepsilon_2}^{\deleted{\unsure{\ell_2}}}\cdots\Delta_{\varepsilon_n}^{\deleted{\unsure{\ell_n}}}&\operatorname{long}x&=n\unsure{.}\\
y&=\Delta_{\varepsilon_1}^{\ell_1}\Delta_{\varepsilon_2}^{\ell_2}\cdots\Delta_{\varepsilon_n}^{\ell_n}&\operatorname{long}y&=\sum_{i\in[1,n]}\ell_i
\end{aligned}
\]

Se donner $g:y\unsure{\text{ --- }|\longrightarrow}z$, i.e. un objet de $y\backslash\underline{\mathrm{Ty}}\deleted{\unsure{/z}}$, revient à se donner une application de l'ens des $\deleted{\unsure{n=}}m=\sum\ell_i$ flèches de transition de $y$ dans $\mathbb N$. En fait, on a
\[
\operatorname{Tr}(y)=\operatorname{Tr}(\Delta_{\varepsilon_1}^{\ell_1})\amalg
\operatorname{Tr}(\Delta_{\varepsilon_2}^{\ell_2})\cdots\amalg\operatorname{Tr}(\Delta_{\varepsilon_n}^{\ell_n})
\]
(\inserted{``une somme'', i.e.} réunion disjointe), donc
\[
\operatorname{Hom}_{\mathrm{Ens}}\unsure{(\operatorname{Tr}y))|(\operatorname{Tr}y)}
\simeq\prod_{i\in[1,\unsure{n|m}]}\operatorname{Hom}_{\mathrm{Ens}}
 \big(\underbrace{\operatorname{Tr}(\Delta_{\varepsilon_i}^{\ell_i})}_{[1,\ell_i]},\mathbb N\big)
\]
donc la donnée de $g$ équivaut à la donnée d'un système de $n$ applications
\[
\left\{\begin{aligned}
\pi_i:[1,\ell_i]&\longrightarrow\mathbb N\qquad(\text{pondération de }\Delta_{\varepsilon_i}^{\ell_i})\\
&\text{pour}\quad i\in[1,n]
\end{aligned}\right.
\]

Mais pour $i\in[1,n]$, considérons \unsure{l'|la} \deleted{\unsure{chaîne}}\inserted{\unsure{élément de|chaîne de|chemin de}} \deleted{flèche} $\operatorname{Nerf}_2$
\diagram{\[
\begin{tikzcd}[ampersand replacement=\&,column sep=normal]
\Delta_{\varepsilon_i}\arrow[r,"{f_{\varepsilon_i,\unsure{\ell_i|c_i}}}"] \&
\Delta_{\varepsilon_i}^{\ell_i}\arrow[r,"{g_{\unsure{\varepsilon_i|\deleted{g}},\pi_i}}"] \&
\underbrace{\Delta_{\varepsilon_i}^{\pi_i(1)}\Delta_{\varepsilon_i}^{\pi_i(\unsure{1|2})}\cdots
\Delta_{\varepsilon_i}^{\pi_i(\ell_i)}}_{\Delta_{\varepsilon_i}^{\sum_j\pi_i(j)\,\overset{\mathrm{déf}}{=}\operatorname{long}(\pi_i)}}
\end{tikzcd}
\tag*{(1)}
\]}
